\section{Quotient Geometry, Exact Dynamics, and Perturbation Analysis}
\label{app:geometry}

This appendix develops the geometry used by the convergence and generalization
results. We follow the organization of a complete manifold analysis: specify the
space and metric, identify the gradient and Hessian, prove the properties of the
implemented retraction, quantify numerical errors, and only then analyze the
iteration. The construction uses standard quotient-manifold principles
\cite{absil2008optimization,boumal2023introduction,mishra2014fixedrank}; the
specialization to balanced two-factor updates and the qualifications about
adaptive state are essential here. All statements concern the regular rank-$r$
stratum. Rank-deficient initialization and rank-changing trajectories require
a separate analysis and are not covered by referring to this stratum by
continuity.

\subsection{Notation and the regular domain}
For matrices, $\langle C,D\rangle=\operatorname{tr}(C^\top D)$ and
$\|C\|=\|C\|_F$. The spectral norm is $\|C\|_2$; $\sigma_r(X)$ denotes
the smallest represented singular value of a rank-$r$ matrix. A factor pair
$p=(B,A)$ has norm $\|p\|^2=\|B\|^2+\|A\|^2$. Let
\begin{align}
 \mathcal P&=\mathbb R_*^{m\times r}\times\mathbb R_*^{r\times d},&
 \pi(B,A)&=BA,\\
 \mathcal M_r&=\{X\in\mathbb R^{m\times d}:\operatorname{rank}X=r\},&
 \mathcal S&=\{(B,A)\in\mathcal P:B^\top B=AA^\top\}.
\end{align}
The subscript $*$ means full column or row rank as appropriate. On a spectral
band, we write
\begin{equation}
 \mathcal K_{a,b}=\{X\in\mathcal M_r:a\leq\sigma_r(X)\leq
 \sigma_1(X)\leq b\},\qquad 0<a\leq b<\infty.
 \label{eq:spectral-band}
\end{equation}
We assume iterates remain in the indicated domain when a theorem uses this
band. Balancing preserves the singular values of $X$ and therefore does not
enforce the lower bound $a$ or the upper bound $b$. Constants that depend on
$a^{-1}$ cannot be declared uniform at a rank boundary.

\begin{theorem}[The quotient and its dimension]
\label{thm:quotient-space}
The action $(B,A)\mapsto(BQ,Q^{-1}A)$ of $\mathrm{GL}(r)$ is free and
proper on $\mathcal P$. Its fibers are exactly the fibers of $\pi$.
Consequently $\mathcal P/\mathrm{GL}(r)$ is diffeomorphic to $\mathcal M_r$
and has dimension $r(m+d-r)$. Moreover,
\begin{equation}
 D\pi_p[\xi]=\xi_BA+B\xi_A,\qquad
 \ker D\pi_p=\mathcal V_p
 =\{(B\Omega,-\Omega A):\Omega\in\mathbb R^{r\times r}\}.
 \label{eq:kernel-exact}
\end{equation}
\end{theorem}
\begin{proof}
If $BQ=B$, the left inverse of $B$ gives $Q=I$, proving freeness.
Suppose $BA=B'A'$ with all factors full rank. Equality of column spaces gives
$B'=BQ$ for a unique invertible $Q$, and applying the left inverse of $B$
gives $A'=Q^{-1}A$. This proves the fiber statement. To check properness,
consider convergent sequences $p_k\to p$ and $\Phi_{Q_k}(p_k)\to p'$ with
$p,p'\in\mathcal P$. The identity $Q_k=B_k^\dagger B'_k$ implies
$Q_k\to B^\dagger B'$. The corresponding identity for $Q_k^{-1}$ shows
that the limit is invertible. Thus the group parameter has a convergent
subsequence in $\mathrm{GL}(r)$, the sequential properness criterion in
these finite-dimensional spaces.

Local smooth sections can also be constructed explicitly. Choose row and
column index sets $I,J$ with $X_{I,J}$ invertible and put
$B=X_{:,J}$, $A=X_{I,J}^{-1}X_{I,:}$. The rank condition implies $BA=X$.
These formulas are smooth in the open chart where the selected minor is
invertible. They identify the smooth quotient with $\mathcal M_r$, not merely
with a local approximation of that manifold. Differentiating the product gives
the first identity in Eq.~\eqref{eq:kernel-exact}. If $D\pi_p[\xi]=0$,
then $(I-BB^\dagger)\xi_BA=0$. Right invertibility of $A$ implies
$\xi_B=B\Omega$, and left invertibility of $B$ gives $\xi_A=-\Omega A$.
The kernel has dimension $r^2$, so rank--nullity gives the asserted dimension.
\end{proof}

\subsection{Balanced representatives and an actual quotient metric}
The Euclidean metric on $\mathcal P$ is not invariant under a general
$\mathrm{GL}(r)$ transformation. An arbitrary Euclidean orthogonal
complement to the $\mathrm{GL}(r)$ orbit is therefore not automatically
a principal connection with the claimed equivariance. We avoid that problem
by defining the Riemannian metric on $\mathcal S/\mathrm O(r)$.

\begin{theorem}[Balancing, residual symmetry, and optimal factor size]
\label{thm:balanced-optimal}
If $X=U\Sigma V^\top$ has rank $r$, every balanced factorization of $X$
has the form
\begin{equation}
 B=U\Sigma^{1/2}R,\qquad A=R^\top\Sigma^{1/2}V^\top,
 \qquad R\in\mathrm O(r).
 \label{eq:balanced-orbit}
\end{equation}
Repeated singular values do not change this residual orthogonal orbit.
For all factorizations $BA=X$,
\begin{align}
 \inf_{BA=X}\tfrac12(\|B\|^2+\|A\|^2)&=\|X\|_*,\label{eq:nuclear-opt}\\
 \inf_{BA=X}\max\{\|B\|_2^2,\|A\|_2^2\}&=\|X\|_2,\label{eq:operator-opt}\\
 \inf_{BA=X}\sqrt{\kappa(B)\kappa(A)}&=\sqrt{\kappa(X)}.
 \label{eq:condition-opt}
\end{align}
Balanced factors attain all three infima. At a balanced pair,
$\|B\|_2=\|A\|_2=\sqrt{\sigma_1(X)}$ and
$\sigma_r(B)=\sigma_r(A)=\sqrt{\sigma_r(X)}$.
\end{theorem}
\begin{proof}
The factors in Eq.~\eqref{eq:balanced-orbit} are balanced and have product
$X$. To prove uniqueness up to $\mathrm O(r)$, let a second balanced pair
be $(BQ,Q^{-1}A)$ and put $H=B^\top B=AA^\top\succ0$.
Its balance equation is $Q^\top H Q=Q^{-1}HQ^{-\top}$.
For $P=QQ^\top\succ0$, multiplication on the left by $Q$ and on the right
by $Q^\top$ yields $PHP=H$. Hence
$(H^{1/2}PH^{1/2})^2=H^2$. Uniqueness of the positive definite square root
gives $P=I$, so $Q$ is orthogonal. No simple-spectrum assumption is used.

The dual characterization of the nuclear norm and Cauchy--Schwarz give
$\|BA\|_*\leq\|B\|\|A\|\leq(\|B\|^2+\|A\|^2)/2$.
The balanced SVD factors attain equality. Similarly,
$\|X\|_2\leq\|B\|_2\|A\|_2\leq\max\{\|B\|_2^2,\|A\|_2^2\}$,
and balanced factors attain $\sigma_1(X)$. Finally, QR decompositions
reduce $BA$ to an invertible $r\times r$ product, giving
$\kappa(X)\leq\kappa(B)\kappa(A)$. Equation~\eqref{eq:balanced-orbit}
attains equality in this condition-number inequality. Taking square roots
proves Eq.~\eqref{eq:condition-opt}.
\end{proof}

These are orbit-optimality statements at a fixed product. They do not imply
that the represented product becomes better conditioned, that the loss is
smaller after balancing, or that every unbalanced factorization is strictly
worse under every possible conditioning statistic.

\begin{lemma}[Balanced tangent and horizontal spaces]
\label{lem:balanced-horizontal}
At $p=(B,A)\in\mathcal S$, define
$C_p(\xi)=B^\top\xi_B-\xi_AA^\top$. Then
\begin{align}
 T_p\mathcal S&=\{\xi:\operatorname{sym}C_p(\xi)=0\},\\
 \mathcal V^{O}_p&=\{(B\Omega,-\Omega A):\Omega^\top=-\Omega\},\\
 \mathcal H_p&=T_p\mathcal S\cap(\mathcal V^O_p)^\perp
             =\{\xi:C_p(\xi)=0\}.
 \label{eq:true-horizontal}
\end{align}
The Euclidean metric restricted to $\mathcal S$ descends to a smooth
Riemannian metric on $\mathcal S/\mathrm O(r)\simeq\mathcal M_r$.
\end{lemma}
\begin{proof}
Differentiating $B^\top B-AA^\top=0$ gives
$C_p(\xi)+C_p(\xi)^\top=0$. Orthogonality to every residual vertical
vector with skew-symmetric $\Omega$ requires
$\operatorname{skew}C_p(\xi)=0$. The two conditions together give
$C_p(\xi)=0$. The derivative of the balance constraint is onto the space
of symmetric matrices: applying a symmetric vertical coefficient leads to
the invertible Sylvester operator $H\Omega+\Omega H$. Thus $\mathcal S$
is a smooth submanifold. The residual action is free, isometric, and proper
because $\mathrm O(r)$ is compact. The standard Riemannian quotient
construction therefore applies.
\end{proof}

\begin{proposition}[Sylvester projection and minimum-norm lift]
\label{prop:minimum-lift}
For any factor direction $u$, let $\Omega$ solve
\begin{equation}
 (B^\top B)\Omega+\Omega(AA^\top)=B^\top u_B-u_AA^\top.
 \label{eq:exact-sylvester}
\end{equation}
Then $\Pi_pu=(u_B-B\Omega,u_A+\Omega A)$ is the Euclidean orthogonal
projection onto $\ker C_p$. It preserves the product differential:
$D\pi_p[\Pi_pu]=D\pi_p[u]$. At a balanced pair this is the horizontal
lift of that product tangent with the smallest factor norm. For any tangent
$Z\in T_X\mathcal M_r$,
\begin{equation}
 \|\operatorname{lift}_p Z\|^2
 =\min_{D\pi_p[\xi]=Z}(\|\xi_B\|^2+\|\xi_A\|^2).
 \label{eq:minimum-lift}
\end{equation}
\end{proposition}
\begin{proof}
Let $V_p\Omega=(B\Omega,-\Omega A)$. Its adjoint is $C_p$, so
$C_pV_p\Omega=(B^\top B)\Omega+\Omega(AA^\top)$.
All eigenvalues of this Sylvester operator are positive sums of Gram
eigenvalues, proving invertibility. Therefore
$\Pi_p=I-V_p(C_pV_p)^{-1}C_p$ is the orthogonal projector onto
$\ker C_p=(\operatorname{im}V_p)^\perp$. Theorem~\ref{thm:quotient-space}
identifies $\operatorname{im}V_p$ with $\ker D\pi_p$, which proves
product-differential preservation. All lifts of $Z$ differ by a vector in
this kernel. Pythagoras then proves Eq.~\eqref{eq:minimum-lift} and
uniqueness of the minimizer. Lemma~\ref{lem:balanced-horizontal} identifies
that minimizer with a tangent to $\mathcal S$ at balanced pairs.
\end{proof}

At an unbalanced pair, Eq.~\eqref{eq:exact-sylvester} still defines a
Euclidean projection, but it is not the same as using the averaged Gram
matrix on both sides of the equation. Nor is this unbalanced Euclidean
projection a $\mathrm{GL}(r)$-equivariant quotient connection. These
distinctions matter when balancing is performed only every $k>1$ steps.

\subsection{An explicit product-space metric and preconditioner}
Let $X=U\Sigma V^\top$ and let $P_U=UU^\top$, $P_V=VV^\top$.
The tangent space consists of matrices whose bottom-right normal block
vanishes:
\begin{equation}
 T_X\mathcal M_r=\{Z:(I-P_U)Z(I-P_V)=0\}.
\end{equation}
Define the self-adjoint ambient operator
\begin{equation}
 \mathcal K_X(G)=U\Sigma U^\top G+GV\Sigma V^\top.
 \label{eq:product-preconditioner}
\end{equation}
It vanishes on normal blocks and is positive definite on the tangent space.

\begin{theorem}[Metric, gradient, and spectral equivalence]
\label{thm:metric-explicit}
The balanced quotient metric has the product-space formula
\begin{equation}
 g_X(Z,Y)=\langle Z,\mathcal K_X^{\dagger}Y\rangle,
 \qquad Z,Y\in T_X\mathcal M_r.
 \label{eq:explicit-metric}
\end{equation}
If $G=\nabla f(X)$, then
\begin{align}
 \operatorname{lift}_p\operatorname{grad}_g f(X)&=(GA^\top,B^\top G),\\
 \operatorname{grad}_g f(X)&=\mathcal K_X(G),\label{eq:true-gradient}\\
 \|\operatorname{grad}_g f(X)\|_g^2&=\|GA^\top\|^2+\|B^\top G\|^2.
 \label{eq:true-grad-norm}
\end{align}
On $\mathcal K_{a,b}$,
\begin{align}
 \frac{1}{2b}\|Z\|^2\leq\|Z\|_g^2&\leq\frac1a\|Z\|^2,
 \label{eq:metric-equivalence}\\
 a\|P_{T_X}G\|^2\leq\|\operatorname{grad}_g f(X)\|_g^2
 &\leq2b\|P_{T_X}G\|^2.
 \label{eq:gradient-equivalence}
\end{align}
\end{theorem}
\begin{proof}
Write $L=D\pi_p$ and use the factor Euclidean inner product. Direct
calculation gives $L^*G=(GA^\top,B^\top G)$ and
$LL^*G=BB^\top G+GA^\top A=\mathcal K_X(G)$ at balanced factors.
The minimum-norm lift in Proposition~\ref{prop:minimum-lift} is
$L^*\mathcal K_X^\dagger Z$. Its squared norm is
$\langle Z,\mathcal K_X^\dagger Z\rangle$, proving the metric formula.
Moreover $B^\top GA^\top-B^\top GA^\top=0$, so $L^*G$ is horizontal.
For every tangent $Z$,
$g_X(LL^*G,Z)=\langle G,Z\rangle=Df(X)[Z]$, proving the gradient identity.

In orthogonal bases $[U,U_\perp]$, $[V,V_\perp]$, the tangent has three
blocks. On $U_iV_j^\top$, $\mathcal K_X$ acts by
$\sigma_i+\sigma_j$; on $U_iV_\perp^\top$ it acts by $\sigma_i$; and
on $U_\perp V_j^\top$ it acts by $\sigma_j$. Nonexistent blocks when
$r=m$ or $r=d$ can simply be omitted. Every tangent eigenvalue lies in
$[a,2b]$. Applying the spectral theorem to $\mathcal K_X$ and its
inverse proves both equivalences.
\end{proof}

This calculation exposes an actual optimization mechanism. Balanced factor
gradient descent is a product-space gradient method with the
state-dependent preconditioner $\mathcal K_X$; it is not the ambient
update $X-\eta\nabla f(X)$. Removing redundant parameters does not make
this preconditioner uniformly well conditioned when $\sigma_r(X)\to0$.

\begin{corollary}[Raw stochastic gradients are horizontal]
\label{cor:raw-horizontal}
For every individual example or minibatch loss $f_z(BA)$, its factor
gradient satisfies $C_p(\nabla(f_z\circ\pi))=0$ exactly. Thus sampling
product-only losses cannot itself create vertical gradient energy.
Coordinatewise adaptive transformations, stale gradient coordinates,
factor-specific penalties, and improperly transported state can create
vertical update energy.
\end{corollary}
\begin{proof}
The chain-rule identity in Eq.~\eqref{eq:true-gradient} applies to each
$f_z$, before any expectation is taken. Applying $C_p$ cancels its two
terms exactly. A nonlinear coordinatewise transformation need not commute
with this cancellation. No such cancellation is available for an arbitrary
penalty depending separately on $B$ and $A$.
\end{proof}

\subsection{Retraction and exact finite-step expansion}
At a balanced pair and a horizontal lift $\xi$, define
\begin{equation}
 R_X(\xi)=(B+\xi_B)(A+\xi_A).
 \label{eq:factor-retraction}
\end{equation}
Here and below a tangent on the quotient is identified with its horizontal
lift when used as the argument of $R$. Balance the factors after this
product is formed to obtain a representative for the next iteration.

\begin{proposition}[A local second-order quotient retraction]
\label{prop:second-order-retraction}
Equation~\eqref{eq:factor-retraction} is a well-defined local retraction
on the balanced quotient, independent of the residual orthogonal gauge.
It is second order for the metric in Eq.~\eqref{eq:explicit-metric}.
It remains in the regular stratum if
\begin{equation}
 \|\xi_B\|_2<\sigma_r(B),\qquad
 \|\xi_A\|_2<\sigma_r(A).
 \label{eq:rank-step}
\end{equation}
Its exact product expansion is
\begin{equation}
 R_X(-\eta\xi)=X-\eta D\pi_p[\xi]+\eta^2\xi_B\xi_A,
 \qquad \|\xi_B\xi_A\|\leq\tfrac12\|\xi\|^2.
 \label{eq:exact-retraction}
\end{equation}
\end{proposition}
\begin{proof}
Residual orthogonal changes send $\xi$ to $(\xi_BR,R^\top\xi_A)$;
the product in Eq.~\eqref{eq:factor-retraction} is unchanged. At zero
the product equals $X$ and its derivative is $D\pi_p[\xi]$, the
retraction axioms. Weyl's singular-value inequality proves
Eq.~\eqref{eq:rank-step}. Multiplication gives the exact expansion, and
$\|\xi_B\xi_A\|\leq\|\xi_B\|\|\xi_A\|\leq\|\xi\|^2/2$ gives its bound.

For the second-order statement consider the factor line
$q(t)=(B+t\xi_B,A+t\xi_A)$. Since $C_p(\xi)=0$, its balance defect is
$t^2(\xi_B^\top\xi_B-\xi_A\xi_A^\top)$. A smooth local balancing gauge
$Q(t)$ with $Q(0)=I$ can therefore be chosen with $Q'(0)=0$;
the Sylvester invertibility in Lemma~\ref{lem:balanced-horizontal}
justifies the implicit-function step without requiring simple singular
values. The balanced curve $s(t)=\Phi_{Q(t)}q(t)$ has initial velocity
$\xi$ and acceleration $s''(0)=V_p(Q''(0))$. This acceleration is
orthogonal to $\mathcal H_p$. The horizontal part of the covariant
acceleration of a curve in the embedded manifold $\mathcal S$ is its
Euclidean acceleration projected onto $\mathcal H_p$. It is therefore
zero at $t=0$. The Riemannian submersion from $\mathcal S$ to the quotient
preserves this horizontal covariant acceleration
\cite{boumal2023introduction}. Thus the quotient retraction is second order.
\end{proof}

\begin{proposition}[Correct horizontal--vertical finite-step discrepancy]
\label{prop:correct-drift}
Let $u=h+v$, $h=\Pi_pu$, and $v=(B\Omega,-\Omega A)$.
The difference between the product after an unprojected update and that
after a projected update at the same factors is exactly
\begin{align}
 &(B-\eta u_B)(A-\eta u_A)-(B-\eta h_B)(A-\eta h_A)\notag\\
 &\quad=\eta^2\{-h_B\Omega A+B\Omega h_A-B\Omega^2A\}.
 \label{eq:correct-cross-terms}
\end{align}
At balanced factors with $\sigma_1(X)\leq b$, its norm is bounded by
\begin{equation}
 \eta^2\left[\sqrt{2b}\,\|h\|\|\Omega\|+b\|\Omega\|^2\right].
 \label{eq:correct-injection}
\end{equation}
Balancing after either update leaves the corresponding product unchanged.
\end{proposition}
\begin{proof}
In the linear part, $D\pi_p[v]=0$ by Eq.~\eqref{eq:kernel-exact}.
In the quadratic part,
$(h_B+B\Omega)(h_A-\Omega A)-h_Bh_A$
equals the braces in Eq.~\eqref{eq:correct-cross-terms}.
Use $\|B\|_2=\|A\|_2\leq\sqrt b$,
$\|\Omega\|_2\leq\|\Omega\|$, and
$\|h_B\|+\|h_A\|\leq\sqrt2\|h\|$ to prove the bound.
\end{proof}

The mixed terms are of order $\eta^2\|h\|\|\Omega\|$, not
$O(\eta^3)$. A bound retaining only $b\eta^2\|\Omega\|^2$ applies to
a purely vertical step, not to the general optimizer update. Furthermore,
vertical first-order transparency holds before balancing as well as after
balancing. Balancing selects a controlled representative for subsequent
steps; it does not cause the algebraic cancellation.

\subsection{Numerical connection errors and transport}
\begin{proposition}[Damping and inexact Sylvester solves]
\label{prop:connection-error}
At a balanced pair $H=B^\top B=AA^\top\succeq aI$, let $\Omega_0$
solve $H\Omega_0+\Omega_0H=C$. Suppose $\widehat\Omega$ satisfies
\begin{equation}
 (H+\varepsilon I)\widehat\Omega+
 \widehat\Omega(H+\varepsilon I)=C+E.
\end{equation}
Then
\begin{align}
 \|\Omega_\varepsilon-\Omega_0\|
 &\leq\frac{\varepsilon}{2a(a+\varepsilon)}\|C\|,
 \label{eq:damping-error}\\
 \|\widehat\Omega-\Omega_\varepsilon\|
 &\leq\frac{\|E\|}{2(a+\varepsilon)},\\
 \|\widehat h-h_0\|
 &\leq\sqrt{2b}\left\{
 \frac{\varepsilon\|C\|}{2a(a+\varepsilon)}+
 \frac{\|E\|}{2(a+\varepsilon)}\right\},
 \label{eq:projector-error}\\
 C_p(\widehat h)&=2\varepsilon\widehat\Omega-E.
 \label{eq:damping-residual-exact}
\end{align}
Here $\widehat h=u-V_p\widehat\Omega$, $h_0=\Pi_pu$ and
$\sigma_1(X)\leq b$. If each of the two Gram matrices is additionally
perturbed by $\Delta H_i$ with spectral norm at most $\delta<a$, the
inverse Sylvester norm is at most $[2(a-\delta+\varepsilon)]^{-1}$
and the same perturbation argument applies with the induced right-hand
error $-\Delta H_1\Omega-\Omega\Delta H_2$.
\end{proposition}
\begin{proof}
Diagonalize $H$. The exact and damped inverse coefficients are
$(\lambda_i+\lambda_j)^{-1}$ and
$(\lambda_i+\lambda_j+2\varepsilon)^{-1}$. Their difference has magnitude
at most $2\varepsilon/[2a(2a+2\varepsilon)]$, proving
Eq.~\eqref{eq:damping-error}. The smallest damped denominator is
$2(a+\varepsilon)$, proving the solve-error bound. The fundamental
field satisfies
$\|V_p\Omega\|^2\leq2b\|\Omega\|^2$; applying this inequality to the
coefficient error proves Eq.~\eqref{eq:projector-error}. Substituting the
inexact solve into $C_p(u-V_p\widehat\Omega)$ proves the exact residual
identity. The last assertion follows from Weyl's inequality and subtracting
the perturbed and unperturbed Sylvester equations.
\end{proof}

Thus a fixed residual threshold that ignores damping can reject a correct
implementation. Damping controls inversion near a rank boundary but does
not restore a full-rank quotient there. In particular, the error relative
to the exact undamped projector can still become large as $a\to0$.

\begin{proposition}[Coordinate transport and state compatibility]
\label{prop:transport}
If $(B',A')=(BQ,Q^{-1}A)$, the action on a factor tangent is
\begin{equation}
 D\Phi_Q(\xi_B,\xi_A)=(\xi_BQ,Q^{-1}\xi_A),
 \quad D\pi_{p'}D\Phi_Q\xi=D\pi_p\xi.
 \label{eq:coordinate-transport}
\end{equation}
Its factor norm is at most
$\max\{\|Q\|_2,\|Q^{-1}\|_2\}\|\xi\|$. At a change between balanced
representatives, $Q\in\mathrm O(r)$ and it is an isometry that commutes
with horizontal projection. Between different quotient points one may
use a smooth local transport $\mathcal T_{x\to y}$; this is distinct
from the same-fiber coordinate change in Eq.~\eqref{eq:coordinate-transport}.
On a compact regular neighborhood such a transport satisfies
\begin{equation}
 \|\mathcal T_{x\to y}\xi-P_{x\to y}\xi\|_y
 \leq C_T d_g(x,y)\|\xi\|_x
 \label{eq:transport-local-error}
\end{equation}
whenever it agrees with the identity on the diagonal and is continuously
differentiable there. $P$ denotes parallel transport on a short unique
geodesic in this neighborhood.
\end{proposition}
\begin{proof}
Differentiate the two factors under the fixed group action and multiply
out the product differential. Submultiplicativity proves the norm bound.
For orthogonal $Q$, Frobenius norms and the Sylvester equation are invariant
under conjugation, proving the balanced assertions. For the last statement,
write both transport operators in a fixed smooth frame on a precompact
chart. Their difference is zero at $y=x$ and has bounded first derivative
on the chart. The mean-value bound gives Eq.~\eqref{eq:transport-local-error}.
\end{proof}

An elementwise second moment is not a tangent vector. Transporting a first
moment using Eq.~\eqref{eq:coordinate-transport} while leaving elementwise
second moments unchanged does not define an invariant Adam algorithm.
Even an orthogonal rotation generally does not commute with elementwise
division by a square root. The convergence results below therefore identify
the analyzed momentum method separately from the archived adaptive heuristic.

\subsection{Landscape equivalence and what it does not prove}
\begin{theorem}[Critical points, Hessian inertia, and gauge nullity]
\label{thm:landscape}
Let $f$ be twice continuously differentiable near $\mathcal M_r$.
A point of $\mathcal P/\mathrm{GL}(r)$ is critical if and only if its
product is critical for $f|_{\mathcal M_r}$. At a critical point the
Hessian bilinear forms in any two smooth quotient/product charts are
related by congruence. They have identical numbers of positive, negative,
and zero eigenvalues. If the product critical point is nondegenerate,
the factor Hessian of $F=f\circ\pi$ has exactly $r^2$ zero directions,
the gauge kernel, and the same positive and negative inertia as the
quotient Hessian.
\end{theorem}
\begin{proof}
The quotient-to-product map is a diffeomorphism by
Theorem~\ref{thm:quotient-space}; its differential is an isomorphism.
The chain rule proves the critical-point equivalence. In chart coordinates,
the second derivative of a composition has a Hessian congruence term and
a term multiplying the first derivative by the second derivative of the
chart map. The latter term vanishes at a critical point. Sylvester's law
of inertia proves the quotient statement.

For the factor formulation, choose local coordinates consisting of a
quotient chart and a group chart. The function is independent of the group
coordinates, so its coordinate Hessian has a quotient block and a zero
$r^2$-dimensional group block. At a critical point, conversion back to
factor coordinates is again a congruence. Nondegeneracy of the quotient
block proves exact nullity and identifies the kernel with the orbit tangent.
\end{proof}

Away from critical points, the factor Hessian is explicitly
\begin{align}
 D^2F(p)[\xi,\zeta]
 &=\langle D\pi_p\xi,\nabla^2f(X)[D\pi_p\zeta]\rangle\notag\\
 &\quad+\langle\nabla f(X),\xi_B\zeta_A+\zeta_B\xi_A\rangle.
 \label{eq:factor-hessian-exact}
\end{align}
Thus simply deleting gauge eigenvectors from an ambient Hessian away from
stationarity need not produce the quotient Hessian. At criticality,
inertia is metric independent but eigenvalue magnitudes are not. A
qualitative strict-saddle property transfers locally; a quantitative global
escape time needs uniform conditioning, regularity, and a specified
perturbation/update law. We do not infer an Adam or momentum saddle-escape
theorem solely from the removal of gauge nullity.

\subsection{Trajectory comparison requires an update-map assumption}
\begin{theorem}[Perturbed quotient trajectories]
\label{thm:trajectory}
Suppose two product or joint-state recursions, driven by the same data and
randomness, obey
\begin{equation}
 d_{t+1}\leq(1+L_U\eta_t)d_t+\eta_t e_t+\eta_t^2 q_t,
 \qquad d_0=0.
 \label{eq:trajectory-recursion}
\end{equation}
Here $e_t$ bounds first-order horizontal/state disagreement and $q_t$
bounds finite-step injections such as Eq.~\eqref{eq:correct-injection}.
Then
\begin{equation}
 d_T\leq\sum_{t<T}(\eta_te_t+\eta_t^2q_t)
 \prod_{k=t+1}^{T-1}(1+L_U\eta_k).
 \label{eq:trajectory-general}
\end{equation}
If instead the common state map is contractive with factor
$1-\mu_U\eta$, $\mu_U>0$, at a constant step with
$0\leq1-\mu_U\eta<1$, and $e_t\leq e$, $q_t\leq q$, then
\begin{equation}
 d_T\leq(1-\mu_U\eta)^T d_0+
 \frac{e+\eta q}{\mu_U}\{1-(1-\mu_U\eta)^T\}.
 \label{eq:trajectory-contract}
\end{equation}
\end{theorem}
\begin{proof}
Repeated substitution in Eq.~\eqref{eq:trajectory-recursion} gives the
product weights in Eq.~\eqref{eq:trajectory-general}. Under contraction,
substitution gives a geometric sum with ratio $1-\mu_U\eta$; its value is
the second term in Eq.~\eqref{eq:trajectory-contract}. These are statements
about the actual update map, including optimizer state when present.
\end{proof}

\begin{proposition}[Strong convexity in $X$ alone is insufficient]
\label{prop:convex-counterexample}
Even a strongly convex product objective need not give a contractive
balanced-factor gradient update on the whole positive rank-one stratum.
For $m=d=r=1$, $x>0$, and $f(x)=\tfrac12(x-c)^2$, $c>0$, balanced
gradient flow is $\dot x=-2x(x-c)$. Its explicit factor-gradient step is
\begin{equation}
 U_\eta(x)=x\{1-\eta(x-c)\}^2.
\end{equation}
For every sufficiently small $\eta>0$, $U_\eta'(x)>1$ at sufficiently
small positive $x<c/2$. Thus $f''=1$ does not imply a product-map
contraction factor $1-\eta$.
\end{proposition}
\begin{proof}
At $B=A=\sqrt x$, both gradients equal $(x-c)\sqrt x$, proving the
update formula. Differentiation gives
$U_\eta'(x)=\{1-\eta(x-c)\}\{1-\eta(3x-c)\}$.
Its limit as $x\downarrow0$ is $(1+\eta c)^2>1$.
\end{proof}

In addition, arbitrary adaptive-state differences can make $e_t$ nonzero.
Neither a summed $O(T\eta^2)$ product-disagreement bound nor a
time-independent $O(\eta)$ ceiling follows merely from balancing and
strong convexity of $f$. The valid conditional version is
Theorem~\ref{thm:trajectory}, with the update-map assumptions stated.

\subsection{Joint geometry for deep factor chains}
For $X=W_L\cdots W_1$ with full-rank compatible factors, let
$\Omega_1,\ldots,\Omega_{L-1}$ denote the interface Lie-algebra elements.
The joint vertical operator is
\begin{align}
 (\mathcal V_W\Omega)_1&=-\Omega_1W_1,\\
 (\mathcal V_W\Omega)_i&=W_i\Omega_{i-1}-\Omega_iW_i,
                      &&1<i<L,\\
 (\mathcal V_W\Omega)_L&=W_L\Omega_{L-1}.
 \label{eq:joint-vertical}
\end{align}
Its adjoint has interface blocks
\begin{equation}
 (\mathcal V_W^*u)_i=W_{i+1}^\top u_{i+1}-u_iW_i^\top.
 \label{eq:joint-adjoint}
\end{equation}

\begin{proposition}[Joint projection and chain balance]
\label{prop:joint-chain}
The minimum-norm removal of all infinitesimal gauge directions solves
\begin{equation}
 (\mathcal V_W^*\mathcal V_W)\Omega=\mathcal V_W^*u,
 \qquad u^{H}=u-\mathcal V_W\Omega.
 \label{eq:joint-normal}
\end{equation}
The normal operator is self-adjoint and positive definite in the regular
chain model with full-row-rank $W_1$, invertible internal factors and
full-column-rank $W_L$. Neighboring interface blocks are coupled.
For $X=U\Sigma V^\top$, a jointly balanced representative is
\begin{equation}
 W_1=\Sigma^{1/L}V^\top,\quad
 W_i=\Sigma^{1/L}\ (1<i<L),\quad
 W_L=U\Sigma^{1/L}.
 \label{eq:joint-chain-balance}
\end{equation}
Its product is $X$ and $W_{i+1}^\top W_{i+1}=W_iW_i^\top$ at every
interface. Each factor condition number is $\kappa(X)^{1/L}$.
\end{proposition}
\begin{proof}
The product differential of Eq.~\eqref{eq:joint-vertical} telescopes to
zero. Expanding the Frobenius inner products gives
Eq.~\eqref{eq:joint-adjoint}. Least squares then gives the normal equations.
If $\mathcal V_W\Omega=0$, its first block and the full row rank of
$W_1$ imply $\Omega_1=0$. Induction through the invertible internal
factors gives every $\Omega_i=0$. Hence the operator is positive definite.
An internal factor belongs to two adjacent interface transformations,
which explains the off-diagonal coupling of the normal equations.
Multiplying Eq.~\eqref{eq:joint-chain-balance} and computing its Grams
proves the remaining assertions.
\end{proof}

A single sweep of independent two-factor projections is generally not the
solution of Eq.~\eqref{eq:joint-normal}. Its error need not scale as
$(L-1)^2$ without further assumptions. The replacement implementation uses
the joint normal equations, records the residual, and treats the older
pairwise depth scans as evidence about that earlier heuristic. This
proposition specifies a geometric correction; it does not predict that
the corrected optimizer must outperform the heuristic on task metrics.
